% Unit 6 self-study -- "I do / You do" ladder for CED 6.1-6.9.
% Ten ladders, forty questions.
%
% Disjoint from u06-frq, which uses the 75.0 g metal calorimetry, the
% NO/NO2 Hess data, the barium hydroxide endothermic demo, a 25.0 g
% ice-to-steam calculation and the two-block thermal-equilibrium question.
% This document uses a neutralisation calorimetry, C/CO/CO2 Hess data,
% H2 + Cl2 bond enthalpies and a 40.0 g phase-change calculation.
\documentclass{apchem-worksheet}

\apcunit{6}
\apcunittitle{Thermochemistry}
\apcdoctitle{Self-Study \textbullet\ CED 6.1--6.9, I~do / You~do}
\apcsubtitle{Ten ladders \textbullet\ four YOUR TURN questions each \textbullet\ work alone, check after all four}

\begin{document}
\apctitleblock

\begin{apcnote}[How to use these notes --- read this first]
Each skill is a \textbf{ladder}: a worked example, then \textbf{four}
YOUR TURN questions --- same skill, new numbers, no help. Work all four
before checking anything.

\textbf{The habit that decides this unit: track the sign, every time.}
Energy leaving the system is negative; energy entering is positive. Most
lost marks here are not arithmetic --- they are a correct number with the
wrong sign, or a $\Delta T$ taken for the wrong substance.

\textbf{Two constants used throughout.}
$c(\ce{H2O}) = \SI{4.18}{\joule\per\gram\per\celsius}$,
$\Delta H_{\text{fus}} = \SI{334}{\joule\per\gram}$ and
$\Delta H_{\text{vap}} = \SI{2260}{\joule\per\gram}$.

\textbf{One scope note.} Topics \ced{6.6} and \ced{6.9} carry a CED
exclusion covering the formal enthalpy-versus-internal-energy distinction
and state-function theory. Hess's law itself is fully examinable; the
thermodynamic formalism behind it is not.
\end{apcnote}

% =====================================================================
\section*{\sffamily Ladder 1 \textbullet\ Endothermic and exothermic}
\ced{6.1}

Define the \emph{system} (the reaction) and the \emph{surroundings}
(everything else, including the beaker and the thermometer). Then ask
which way energy flowed.

\begin{workedexample}[I do: reading a temperature change backwards]
A beaker becomes cold when two solids are mixed. State the sign of
$\Delta H$ and explain the flow.

\textbf{The surroundings got colder}, so the surroundings \emph{lost}
energy. That energy went into the system.

\textbf{The system absorbed energy}, so the process is
\textbf{endothermic} and $\Delta H$ is \textbf{positive}.

\textbf{The reversal that trips people:} what you feel is the
\emph{surroundings}, but the sign refers to the \emph{system}. A cold
beaker means a positive $\Delta H$; a hot beaker means a negative one.

\textbf{What never to write:} ``the reaction gives off cold''. Cold is not
a substance. Energy flowed one way only --- into the system.
\end{workedexample}

\begin{yourturn}[YOUR TURN 1 --- four questions]
\begin{enumerate}[leftmargin=*,itemsep=0.5em]
  \item Combustion warms its surroundings. Sign of $\Delta H$?
        \workspace[1.4cm]
  \item \ce{NH4NO3} dissolving makes the beaker cold. Sign?
        \workspace[1.4cm]
  \item In an exothermic reaction, do the surroundings gain or lose
        energy? \workspace[1.4cm]
  \item Melting ice is endothermic. Explain in terms of the system.
        \workspace[1.8cm]
\end{enumerate}
\selfcheck{(a) negative \quad (b) positive \quad (c) gain \quad (d) the
ice absorbs energy to overcome attractions}
\keyonly{\begin{apcanswer}(a) \textbf{Negative} --- the system released
energy to the surroundings.
(b) \textbf{Positive} --- the system absorbed energy from the
surroundings, which therefore cooled.
(c) They \textbf{gain} energy, which is why they warm up.
(d) The ice (the system) must \textbf{absorb} energy to break the hydrogen
bonds holding its lattice together. That energy comes from the
surroundings, which is why your hand feels cold holding it: your hand is
the surroundings, losing energy.\end{apcanswer}}
\end{yourturn}

% =====================================================================
\section*{\sffamily Ladder 2 \textbullet\ Energy diagrams}
\ced{6.2}

The vertical axis is energy. Where the products sit relative to the
reactants gives the sign of $\Delta H$; the hump between them is the
activation energy, which says nothing about $\Delta H$.

\begin{workedexample}[I do: reading both quantities off one picture]
A diagram shows reactants at \SI{60}{\kilo\joule}, a peak at
\SI{150}{\kilo\joule} and products at \SI{25}{\kilo\joule}. Find $E_a$
and $\Delta H$, and classify the reaction.

\textbf{$E_a$ (forward):} $150 - 60 = \mathbf{90~kJ}$.

\textbf{$\Delta H$:} products minus reactants, $25 - 60 =
\mathbf{-35~kJ}$. Negative, so \textbf{exothermic}.

\textbf{The visual signature:} products \emph{below} reactants means
exothermic; above means endothermic. Read the two \emph{ends}, never the
peak, for $\Delta H$.

\textbf{Why the peak is irrelevant to $\Delta H$:} a catalyst lowers the
peak without moving either end, so it changes $E_a$ and leaves $\Delta H$
untouched. The two quantities are independent.
\end{workedexample}

\begin{yourturn}[YOUR TURN 2 --- four questions]
\begin{enumerate}[leftmargin=*,itemsep=0.5em]
  \item Reactants \SI{40}{\kilo\joule}, peak \SI{130}{\kilo\joule},
        products \SI{95}{\kilo\joule}. $E_a$? \workspace[1.4cm]
  \item $\Delta H$ for that reaction, and its classification?
        \workspace[1.6cm]
  \item On an exothermic diagram, where do the products sit?
        \workspace[1.4cm]
  \item Does a large $E_a$ mean a large $\Delta H$? Explain.
        \workspace[1.8cm]
\end{enumerate}
\selfcheck{(a) \SI{90}{\kilo\joule} \quad (b) $+55$~kJ, endothermic \quad
(c) below the reactants \quad (d) no --- they are independent}
\keyonly{\begin{apcanswer}(a) $130 - 40 = \mathbf{90}$~kJ.
(b) $95 - 40 = \mathbf{+55}$~kJ, so \textbf{endothermic}.
(c) \textbf{Below} the reactants --- the system ended at lower energy,
having released the difference.
(d) \textbf{No.} $E_a$ is the height of the barrier between reactants and
products; $\Delta H$ is the difference between the two ends. A reaction
can have an enormous barrier and a tiny $\Delta H$, or a small barrier and
a huge $\Delta H$. Adding a catalyst proves the point: it changes $E_a$
and leaves $\Delta H$ exactly as it was.\end{apcanswer}}
\end{yourturn}

% =====================================================================
\section*{\sffamily Ladder 3 \textbullet\ Heat transfer and thermal equilibrium}
\ced{6.3}

Energy flows from hotter to cooler until both reach the same temperature.
Same temperature does not mean same energy.

\begin{workedexample}[I do: what equalizes and what does not]
Two blocks at different temperatures are placed in contact and isolated.
Describe what happens at the particle level, and state what is equal at
equilibrium.

\textbf{At the boundary:} particles of the hotter block have greater
average kinetic energy. They collide with particles of the cooler block
and transfer energy to them.

\textbf{Over time:} the hotter block's average kinetic energy falls, the
cooler block's rises, until the two averages match. Collisions continue,
but with no further \emph{net} transfer --- the flow is now equal in both
directions.

\textbf{What is equal at equilibrium: the temperature}, meaning the
average kinetic energy per particle.

\textbf{What is \emph{not} equal: the total thermal energy.} That depends
on mass and specific heat capacity too. A large block of water at
\SI{25}{\celsius} holds far more energy than a small block of copper at
the same temperature.
\end{workedexample}

\begin{yourturn}[YOUR TURN 3 --- four questions]
\begin{enumerate}[leftmargin=*,itemsep=0.5em]
  \item Which direction does energy flow between objects at different
        temperatures? \workspace[1.4cm]
  \item At thermal equilibrium, what quantity is equal?
        \workspace[1.4cm]
  \item Do the two objects then hold equal thermal energy? Explain.
        \workspace[1.8cm]
  \item Do collisions stop at thermal equilibrium? \workspace[1.6cm]
\end{enumerate}
\selfcheck{(a) hotter to cooler \quad (b) temperature \quad (c) no ---
mass and $c$ differ \quad (d) no --- net transfer stops, not the
collisions}
\keyonly{\begin{apcanswer}(a) From \textbf{hotter to cooler}, always.
(b) The \textbf{temperature} --- the average kinetic energy per particle.
(c) \textbf{No.} Total thermal energy depends on mass and specific heat
capacity as well as temperature. Two objects at the same temperature can
hold vastly different amounts.
(d) \textbf{No.} The particles keep colliding and keep exchanging energy.
What stops is the \emph{net} flow: transfer in each direction becomes
equal, so nothing changes on the macroscopic
scale.\end{apcanswer}}
\end{yourturn}

% =====================================================================
\section*{\sffamily Ladder 4 \textbullet\ Specific heat}
\ced{6.4}

$q = mc\,\Delta T$. Specific heat capacity is the energy needed to raise
one gram by one degree --- a property of the substance, not of the sample.

\begin{workedexample}[I do: energy to warm water]
How much energy warms \SI{125}{\gram} of water from \SI{20.0}{\celsius}
to \SI{45.0}{\celsius}?

\[ q = mc\,\Delta T = (125)(4.18)(25.0) = \mathbf{1.31\times10^{4}~J}
   = \SI{13.1}{\kilo\joule} \]

\textbf{$\Delta T$ is final minus initial}, and its size in
\si{\celsius} equals its size in kelvin --- so either unit works here, as
long as you use a \emph{difference}.

\textbf{Why water's value is so large.} At
\SI{4.18}{\joule\per\gram\per\celsius}, water takes far more energy per
degree than most substances (copper is \num{0.385}). Its hydrogen-bond
network absorbs energy without the molecules speeding up much --- which is
why water moderates climate and makes a good coolant.
\end{workedexample}

\begin{yourturn}[YOUR TURN 4 --- four questions]
\begin{enumerate}[leftmargin=*,itemsep=0.5em]
  \item Energy to warm \SI{50.0}{\gram} of water by
        \SI{30.0}{\celsius}: \workspace[1.6cm]
  \item \SI{2500}{\joule} into \SI{100.0}{\gram} of water. Temperature
        rise? \workspace[1.6cm]
  \item \SI{1200}{\joule} raises a metal
        ($c = \SI{0.900}{\joule\per\gram\per\celsius}$) by
        \SI{15.0}{\celsius}. Mass? \workspace[1.8cm]
  \item Equal masses of water and copper absorb the same energy. Which
        warms more, and why? \workspace[1.8cm]
\end{enumerate}
\selfcheck{(a) \SI{6270}{\joule} \quad (b) \SI{5.98}{\celsius} \quad
(c) \SI{88.9}{\gram} \quad (d) copper --- much smaller $c$}
\keyonly{\begin{apcanswer}(a) $(50.0)(4.18)(30.0) = \mathbf{6270}$~J.
(b) $\Delta T = q/(mc) = 2500/[(100.0)(4.18)] =
\mathbf{5.98}$~\si{\celsius}.
(c) $m = q/(c\,\Delta T) = 1200/[(0.900)(15.0)] = \mathbf{88.9}$~g.
(d) \textbf{Copper}, by roughly a factor of eleven. Its specific heat
capacity (\num{0.385}) is far smaller than water's (\num{4.18}), so the
same energy per gram produces a much larger temperature
change.\end{apcanswer}}
\end{yourturn}

% =====================================================================
\section*{\sffamily Ladder 5 \textbullet\ Calorimetry}
\ced{6.4} \ced{6.6}

In a coffee-cup calorimeter the \emph{solution} absorbs the heat, so
$q$ is computed for the solution --- and the reaction's $\Delta H$ is the
opposite sign, per mole.

\begin{workedexample}[I do: a neutralization]
\SI{50.0}{\milli\litre} of \SI{1.00}{\Molar} \ce{HCl} and
\SI{50.0}{\milli\litre} of \SI{1.00}{\Molar} \ce{NaOH}, both at
\SI{21.0}{\celsius}, are mixed. The temperature peaks at
\SI{28.2}{\celsius}. Find $\Delta H$ per mole of water formed.

\textbf{Mass of solution:} $50.0 + 50.0 = \SI{100.0}{\gram}$ (assuming
the density of water). \emph{Total} volume --- not one of them.

\textbf{Heat absorbed by the solution:}
\[ q = (100.0)(4.18)(7.2) = \SI{3.0e3}{\joule} =
   \SI{3.01}{\kilo\joule} \]

\textbf{Moles of water formed:} $(0.0500)(1.00) = \num{0.0500}$~mol ---
both reactants supply exactly this, so neither is in excess.

\textbf{Per mole, with the sign flipped:}
\[ \Delta H = -\frac{3.01}{0.0500} = \mathbf{-60.2~kJ/mol} \]
Negative because the solution \emph{gained} what the reaction
\emph{released}.

\textbf{Two traps:} using \SI{50.0}{\gram} instead of \SI{100.0}{\gram},
and forgetting to flip the sign.
\end{workedexample}

\begin{yourturn}[YOUR TURN 5 --- four questions]
\begin{enumerate}[leftmargin=*,itemsep=0.5em]
  \item Why is the total mass of solution used, not one reactant's?
        \workspace[1.6cm]
  \item If the same experiment used \SI{100}{\milli\litre} of each, would
        $\Delta T$ change? Would $\Delta H$ per mole?
        \workspace[2.0cm]
  \item Why is $\Delta H$ negative when the temperature rose?
        \workspace[1.6cm]
  \item Give one reason the measured magnitude is usually too small.
        \workspace[1.6cm]
\end{enumerate}
\selfcheck{(a) the whole solution absorbs the heat \quad (b) $\Delta T$
about the same; $\Delta H$ per mole unchanged \quad (c) the system
released it \quad (d) heat lost to the surroundings}
\keyonly{\begin{apcanswer}(a) Because the heat released spreads through
\emph{all} the liquid present. The whole mixed solution is what warms up,
so its whole mass appears in $q = mc\,\Delta T$.
(b) $\Delta T$ would be essentially \textbf{unchanged}: twice the heat is
released, but it warms twice as much solution. $\Delta H$ per mole is also
\textbf{unchanged} --- it is an intensive property of the reaction, not of
how much was run.
(c) Because the sign convention refers to the \textbf{system}. The
solution (surroundings) gained energy, so the reaction released it, making
$\Delta H$ negative.
(d) \textbf{Heat is lost to the surroundings} --- the cup, the
thermometer, the air --- so the measured temperature rise, and hence the
computed $q$, is smaller than the true value. (Accept also: the heat
capacity of the calorimeter itself is neglected.)\end{apcanswer}}
\end{yourturn}

% =====================================================================
\section*{\sffamily Ladder 6 \textbullet\ Energy of phase changes}
\ced{6.5}

A phase change uses $q = m\,\Delta H_{\text{phase}}$ --- mass only, no
$\Delta T$, because the temperature does not change.

\begin{workedexample}[I do: ice to steam in three steps]
Find the energy to convert \SI{40.0}{\gram} of ice at \SI{0}{\celsius}
into steam at \SI{100}{\celsius}.

\textbf{Step 1 --- melt} (no $\Delta T$):
$q = (40.0)(334) = \SI{13.4}{\kilo\joule}$.

\textbf{Step 2 --- warm the liquid} (this one has a $\Delta T$):
$q = (40.0)(4.18)(100) = \SI{16.7}{\kilo\joule}$.

\textbf{Step 3 --- vaporize} (no $\Delta T$):
$q = (40.0)(2260) = \SI{90.4}{\kilo\joule}$.

\[ \text{Total} = 13.4 + 16.7 + 90.4 = \mathbf{120~kJ} \]

\textbf{Notice which step dominates.} Vaporizing takes more than the other
two combined, because melting only loosens the lattice while vaporizing
must separate the molecules completely --- overcoming essentially all the
hydrogen bonding.

\textbf{The structural error to avoid:} putting a $\Delta T$ into a
phase-change step. During melting or boiling the temperature is constant,
so there is no $\Delta T$ to use.
\end{workedexample}

\begin{yourturn}[YOUR TURN 6 --- four questions]
\begin{enumerate}[leftmargin=*,itemsep=0.5em]
  \item Energy to melt \SI{60.0}{\gram} of ice at \SI{0}{\celsius}:
        \workspace[1.6cm]
  \item Energy to vaporize \SI{10.0}{\gram} of water at
        \SI{100}{\celsius}: \workspace[1.6cm]
  \item Why does a phase-change calculation use no $\Delta T$?
        \workspace[1.6cm]
  \item Why is $\Delta H_{\text{vap}}$ so much larger than
        $\Delta H_{\text{fus}}$? \workspace[1.8cm]
\end{enumerate}
\selfcheck{(a) \SI{20.0}{\kilo\joule} \quad (b) \SI{22.6}{\kilo\joule}
\quad (c) temperature is constant during the change \quad
(d) vaporizing separates molecules completely}
\keyonly{\begin{apcanswer}(a) $(60.0)(334) = \SI{20040}{\joule} =
\mathbf{20.0}$~kJ.
(b) $(10.0)(2260) = \SI{22600}{\joule} = \mathbf{22.6}$~kJ.
(c) Because the temperature does not change during a phase transition ---
all the supplied energy goes into overcoming intermolecular attractions
(potential energy) rather than into speeding the particles up.
(d) \textbf{Melting} only loosens the structure: molecules leave fixed
positions but remain in contact, and most attractions survive.
\textbf{Vaporizing} must separate them completely, overcoming essentially
all of the remaining attractions and moving the molecules far apart. For
water that is nearly seven times as much energy per
gram.\end{apcanswer}}
\end{yourturn}

% =====================================================================
\section*{\sffamily Ladder 7 \textbullet\ Hess's law}
\ced{6.9}

Enthalpy changes add. Reverse a reaction and the sign flips; multiply a
reaction and $\Delta H$ multiplies too.

\begin{workedexample}[I do: manipulating to reach a target]
Given
\ce{C(s) + O2(g) -> CO2(g)}, $\Delta H = -393.5$~kJ, and
\ce{CO(g) + 1/2O2(g) -> CO2(g)}, $\Delta H = -283.0$~kJ,
find $\Delta H$ for \ce{C(s) + 1/2O2(g) -> CO(g)}.

\textbf{Target check:} \ce{CO} must end on the \emph{right}, but in the
second equation it is on the left. So reverse the second and flip its
sign:
\[ \ce{CO2(g) -> CO(g) + 1/2O2(g)} \qquad \Delta H = +283.0~\text{kJ} \]

\textbf{Add to the first:}
\ce{C + O2 + CO2 -> CO2 + CO + 1/2O2}. The \ce{CO2} cancels and the
oxygens net to \ce{1/2 O2} on the left.

\[ \Delta H = (-393.5) + (+283.0) = \mathbf{-110.5~kJ} \]

\textbf{Why it works:} enthalpy depends only on the initial and final
states, not the route. Any sequence of steps connecting the same start and
finish must sum to the same total --- which is what lets us obtain a value
for a reaction too difficult to measure directly, as this one is.
\end{workedexample}

\begin{yourturn}[YOUR TURN 7 --- four questions]
\begin{enumerate}[leftmargin=*,itemsep=0.5em]
  \item If \ce{A -> B} has $\Delta H = -50$~kJ, what is it for
        \ce{B -> A}? \workspace[1.4cm]
  \item If \ce{A -> B} has $\Delta H = -50$~kJ, what is it for
        \ce{2A -> 2B}? \workspace[1.4cm]
  \item Why can enthalpy changes be added at all? \workspace[1.6cm]
  \item Why is Hess's law useful in practice? \workspace[1.8cm]
\end{enumerate}
\selfcheck{(a) $+50$~kJ \quad (b) $-100$~kJ \quad (c) enthalpy depends
only on the states \quad (d) it gives values that cannot be measured
directly}
\keyonly{\begin{apcanswer}(a) $\mathbf{+50}$~kJ --- reversing flips the
sign.
(b) $\mathbf{-100}$~kJ --- doubling the amounts doubles the energy.
(c) Because enthalpy depends only on the current state of the system, not
on how it got there. The difference between reactants and products is
therefore fixed, so any path between them sums to the same total.
(d) It supplies $\Delta H$ for reactions that cannot be run cleanly in a
calorimeter. The example above is exactly such a case: burning carbon to
give \emph{only} \ce{CO} is not achievable in practice --- some \ce{CO2}
always forms --- so the value must be obtained
indirectly.\end{apcanswer}}
\end{yourturn}

% =====================================================================
\section*{\sffamily Ladder 8 \textbullet\ Bond enthalpies}
\ced{6.7}

Breaking bonds costs energy; forming them releases it.
\[ \Delta H \approx \sum(\text{bonds broken}) -
   \sum(\text{bonds formed}) \]

\begin{workedexample}[I do: a simple gas-phase reaction]
Estimate $\Delta H$ for \ce{H2(g) + Cl2(g) -> 2HCl(g)} from
\ce{H-H} 436, \ce{Cl-Cl} 243, \ce{H-Cl} 431
\si{\kilo\joule\per\mole}.

\textbf{Broken} (energy in): $436 + 243 = \num{679}$~kJ.

\textbf{Formed} (energy out): $2 \times 431 = \num{862}$~kJ.

\[ \Delta H = 679 - 862 = \mathbf{-183~kJ} \]

Negative, so exothermic --- more energy is released making the new bonds
than was needed to break the old ones. The accepted value is
\SI{-185}{\kilo\joule}, so the estimate is good to about
\SI{1}{\percent}.

\textbf{Why only an estimate:} tabulated bond enthalpies are
\emph{averages} across many different compounds. A particular \ce{H-Cl}
bond in a particular molecule differs slightly from the average. The
method also assumes every species is a gas.

\textbf{Direction of the subtraction:} broken \emph{minus} formed.
Reversing it gives the right magnitude with the wrong sign, which is a
whole-question error.
\end{workedexample}

\begin{yourturn}[YOUR TURN 8 --- four questions]
\begin{enumerate}[leftmargin=*,itemsep=0.5em]
  \item Is breaking a bond endothermic or exothermic?
        \workspace[1.4cm]
  \item If bonds broken total \SI{800}{\kilo\joule} and formed
        \SI{950}{\kilo\joule}, find $\Delta H$. \workspace[1.6cm]
  \item Why is a bond-enthalpy result only approximate?
        \workspace[1.8cm]
  \item Why does the method require all species to be gases?
        \workspace[1.8cm]
\end{enumerate}
\selfcheck{(a) endothermic \quad (b) $-150$~kJ \quad (c) tabulated values
are averages \quad (d) otherwise intermolecular forces contribute too}
\keyonly{\begin{apcanswer}(a) \textbf{Endothermic} --- energy must be
supplied to pull bonded atoms apart. Forming bonds is exothermic.
(b) $800 - 950 = \mathbf{-150}$~kJ, exothermic.
(c) Because tabulated bond enthalpies are \textbf{averages} taken over
many compounds. The true strength of a given bond depends on its
molecular environment, so a total assembled from averages will not match
the exact value.
(d) Because the method accounts only for making and breaking
\emph{covalent bonds within} molecules. If a species were a liquid or
solid, intermolecular forces would also have to be overcome or formed, and
those are not included in the bond-enthalpy
bookkeeping.\end{apcanswer}}
\end{yourturn}

% =====================================================================
\section*{\sffamily Ladder 9 \textbullet\ Enthalpy of formation}
\ced{6.8}

\[ \Delta H^\circ_{\text{rxn}} = \sum \Delta H_f^\circ(\text{products})
   - \sum \Delta H_f^\circ(\text{reactants}) \]
An element in its standard state has $\Delta H_f^\circ = 0$.

\begin{workedexample}[I do: a combustion from formation data]
Find $\Delta H^\circ$ for
\ce{C2H4(g) + 3O2(g) -> 2CO2(g) + 2H2O(l)} using
$\Delta H_f^\circ$: \ce{C2H4} $+52.3$, \ce{CO2} $-393.5$,
\ce{H2O(l)} $-285.8$ \si{\kilo\joule\per\mole}.

\textbf{Products}, with coefficients applied:
$2(-393.5) + 2(-285.8) = -787.0 - 571.6 = \num{-1358.6}$.

\textbf{Reactants:} $(+52.3) + 3(0) = \num{52.3}$. Oxygen is an element in
its standard state, so its value is \textbf{zero} --- it is not omitted, it
is genuinely zero.

\[ \Delta H^\circ = -1358.6 - 52.3 = \mathbf{-1411~kJ} \]

\textbf{Two things to get right:} apply the coefficients (a missed factor
of 2 here changes the answer by hundreds of kilojoules), and subtract in
the right direction --- products minus reactants.
\end{workedexample}

\begin{yourturn}[YOUR TURN 9 --- four questions]
\begin{enumerate}[leftmargin=*,itemsep=0.5em]
  \item What is $\Delta H_f^\circ$ for \ce{N2(g)}? Why?
        \workspace[1.6cm]
  \item What is $\Delta H_f^\circ$ for \ce{O3(g)}, and why is it not
        zero? \workspace[1.8cm]
  \item Products total $-500$, reactants $-120$~kJ. Find
        $\Delta H^\circ$. \workspace[1.4cm]
  \item Why must the coefficients be applied to the tabulated values?
        \workspace[1.8cm]
\end{enumerate}
\selfcheck{(a) zero --- element in its standard state \quad (b) non-zero;
\ce{O2} is the standard state, not \ce{O3} \quad (c) $-380$~kJ \quad
(d) the table is per mole}
\keyonly{\begin{apcanswer}(a) \textbf{Zero}, because \ce{N2}(g) is the
standard state of the element nitrogen, and forming an element from itself
involves no change.
(b) \textbf{Not zero} (it is about $+143$~kJ/mol). Ozone is not the
standard state of oxygen --- \ce{O2} is --- so making \ce{O3} from \ce{O2}
is a real chemical change with a real enthalpy.
(c) $-500 - (-120) = \mathbf{-380}$~kJ.
(d) Because $\Delta H_f^\circ$ is tabulated \textbf{per mole} of
substance. If the equation makes two moles, twice that enthalpy is
involved. Omitting a coefficient is the commonest error in this
calculation.\end{apcanswer}}
\end{yourturn}

% =====================================================================
\section*{\sffamily Ladder 10 \textbullet\ Choosing the right method}
\ced{6.6} \ced{6.7} \ced{6.8} \ced{6.9}

Three routes give $\Delta H$. Knowing which one the data are asking for is
half the skill.

\begin{workedexample}[I do: matching data to method]
State which method to use for each set of data.

\textbf{Given a temperature change in a calorimeter:} use
$q = mc\,\Delta T$, then divide by moles and flip the sign. This is the
only \emph{experimental} route.

\textbf{Given a table of $\Delta H_f^\circ$ values:} use products minus
reactants. Most accurate of the three, and the usual exam default.

\textbf{Given a table of bond enthalpies:} use broken minus formed. Only
an estimate, and only for gases.

\textbf{Given other reactions with their $\Delta H$ values:} use Hess's
law --- reverse and scale until they add to the target.

\textbf{The tell:} the data given \emph{is} the instruction. Temperatures
mean calorimetry; a $\Delta H_f^\circ$ table means the formation formula;
bonds mean the estimate; whole reactions mean Hess.
\end{workedexample}

\begin{yourturn}[YOUR TURN 10 --- four questions]
\begin{enumerate}[leftmargin=*,itemsep=0.5em]
  \item You are given the mass, specific heat and temperature change.
        Which method? \workspace[1.4cm]
  \item You are given two reactions with known $\Delta H$. Which method?
        \workspace[1.4cm]
  \item Which method gives only an estimate, and why?
        \workspace[1.6cm]
  \item Which method is the only truly experimental one?
        \workspace[1.4cm]
\end{enumerate}
\selfcheck{(a) calorimetry \quad (b) Hess's law \quad (c) bond
enthalpies --- averaged values \quad (d) calorimetry}
\keyonly{\begin{apcanswer}(a) \textbf{Calorimetry}: $q = mc\,\Delta T$,
then per mole with the sign reversed.
(b) \textbf{Hess's law} --- reverse and scale them until they sum to the
target reaction.
(c) \textbf{Bond enthalpies}, because the tabulated values are averages
across many compounds rather than values for the specific molecules in
question. It also assumes all species are gaseous.
(d) \textbf{Calorimetry.} The other three all begin from tabulated data
that someone else measured; calorimetry is the measurement
itself.\end{apcanswer}}
\end{yourturn}

% =====================================================================
\section*{\sffamily Where you stand}

Tick a ladder only if all four YOUR TURN questions were right first time.

\renewcommand{\arraystretch}{1.30}
\noindent\begin{tabular}{@{}c p{6.0cm} c >{\raggedright\arraybackslash}p{4.9cm}@{}}
\toprule
\textbf{First try?} & \textbf{Skill} & \textbf{Ladder} &
  \textbf{If not, re-read\ldots}\\
\midrule
$\square$ & endo- vs exothermic & 1 & the sign is the system's \\
$\square$ & energy diagrams & 2 & read the ends, not the peak \\
$\square$ & thermal equilibrium & 3 & equal $T$, not equal energy \\
$\square$ & specific heat & 4 & $q = mc\Delta T$ \\
$\square$ & calorimetry & 5 & total mass; flip the sign \\
$\square$ & phase changes & 6 & mass only, no $\Delta T$ \\
$\square$ & Hess's law & 7 & reverse flips, scaling multiplies \\
$\square$ & bond enthalpies & 8 & broken minus formed \\
$\square$ & enthalpy of formation & 9 & products minus reactants \\
$\square$ & choosing a method & 10 & the data tells you \\
\bottomrule
\end{tabular}

\begin{apcnote}[Scoring yourself honestly]
10/10: move on to the unit worksheets and the free-response set.
7--9: solid --- redo the missed ladders tomorrow from a blank page.
6 or fewer: the recurring root causes here are exactly three ---
(1) a correct magnitude with the wrong \emph{sign}, (2) using one
reactant's mass instead of the total solution mass in calorimetry, and
(3) forgetting to apply the coefficients to tabulated
$\Delta H_f^\circ$ values. Fix those three and most of these ladders fall
together.
\end{apcnote}

\end{document}
